% Part: first-order-logic
% Chapter: axiomatic-deduction
% Section: provability-consistency

\documentclass[../../../include/open-logic-section]{subfiles}

\begin{document}

\iftag{FOL}
      {\olfileid{fol}{axd}{prv}}
      {\olfileid{pl}{axd}{prv}}

\olsection{\usetoken{S}{derivability} lan Konsistensi}

Saiki kita bakal netepake sawetara sipat relasi !!{derivability}.
Sipat kasebut narik kawigaten dhewe-dhewe, nanging saben sipat bakal
duwe peran ing bukti teorema kelengkapan.

\begin{prop}\ollabel{prop:provability-contr}
  Yen $\Gamma \Proves !A$ lan $\Gamma \cup \{!A\}$ ora konsisten,
  mula $\Gamma$ ora konsisten.
\end{prop}

\begin{proof}
  Yen $\Gamma \cup \{!A\}$ ora konsisten, mula $\Gamma \cup \{!A\}
  \Proves \lfalse$. Miturut \olref[ptn]{prop:reflexivity}, $\Gamma
  \Proves !B$ kanggo saben $!B \in \Gamma$. Amarga miturut hipotesis
  kita uga duwe $\Gamma \Proves !A$, mula $\Gamma \Proves !B$ kanggo
  saben $!B \in \Gamma \cup \{!A\}$. Miturut
  \olref[ptn]{prop:transitivity}, $\Gamma \Proves \lfalse$, yaiku
  $\Gamma$ ora konsisten.
\end{proof}

\begin{prop}
\ollabel{prop:prov-incons}
$\Gamma \Proves !A$ yen lan mung yen $\Gamma \cup \{\lnot !A\}$ ora
konsisten.
\end{prop}

\begin{proof}
Kaping pisan, upamana $\Gamma \Proves !A$. Banjur $\Gamma \cup \{\lnot
!A\} \Proves !A$ miturut \olref[ptn]{prop:monotonicity}. $\Gamma \cup
\{\lnot !A\} \Proves \lnot !A$ miturut
\olref[ptn]{prop:reflexivity}. Kita uga duwe $\Proves \lnot !A \lif
(!A \lif \lfalse)$ miturut \olref[prp]{ax:lnot2}. Mula, nganggo rong
aplikasi \olref[ded]{prop:mp}, kita duwe $\Gamma \cup \{\lnot !A\}
\Proves \lfalse$.

Saiki upamana $\Gamma \cup \{\lnot !A\}$ ora konsisten, yaiku $\Gamma
\cup \{\lnot !A\} \Proves \lfalse$. Miturut Teorema Deduksi, $\Gamma
\Proves \lnot !A \lif \lfalse$. $\Gamma \Proves (\lnot !A \lif
\lfalse) \lif \lnot\lnot !A$ miturut \olref[prp]{ax:lfalse2}, mula
$\Gamma \Proves \lnot\lnot !A$ miturut \olref[ded]{prop:mp}. Amarga
$\Gamma \Proves \lnot\lnot !A \lif !A$ (\olref[prp]{ax:dne}), kita
entuk $\Gamma \Proves !A$ kanthi nggunakake
\olref[ded]{prop:mp} maneh.
\end{proof}

\begin{prob}
Buktekna yen $\Gamma \Proves \lnot !A$ yen lan mung yen $\Gamma \cup
\{!A\}$ ora konsisten.
\end{prob}

\begin{prop}\ollabel{prop:explicit-inc}
  Yen $\Gamma \Proves !A$ lan $\lnot !A \in \Gamma$, mula $\Gamma$
  ora konsisten.
\end{prop}

\begin{proof}
  Miturut refleksivitas, negasi A bisa diderivasi saka Gamma.
  Uga, $\Gamma \Proves \lnot !A \lif (!A \lif \lfalse)$ miturut
  \olref[prp]{ax:lnot2}. Mula $\Gamma \Proves \lfalse$ kanthi rong
  aplikasi \olref[ded]{prop:mp}. Cathetan koreksi sumber OLPL-031:
  langkah refleksivitas kang diwajibake dening premis kaanggotaan
  ditulis cetha sadurunge rong aplikasi modus ponens.
\end{proof}


\begin{prop}\ollabel{prop:provability-exhaustive}
  Yen $\Gamma \cup \{!A\}$ lan $\Gamma \cup \{\lnot !A\}$ kalorone
  ora konsisten, mula $\Gamma$ ora konsisten.
\end{prop}

\begin{proof}
Gladhen.
\end{proof}

\tagprob{FOL}
\begin{prob}
  Buktekna \olref[fol][axd][prv]{prop:provability-exhaustive}
\end{prob}
\tagendprob

\tagprob{notFOL}
\begin{prob}
  Buktekna \olref[pl][axd][prv]{prop:provability-exhaustive}
\end{prob}
\tagendprob


\end{document}
